\section{Universal Contracts, Benchmark, and Governance}
\label{sec:ieee-contracts}

The universal claim only matters if the runtime interface, benchmark interface, and
governance interface remain stable across domains. ORIUS enforces that stability
through three public surfaces: a single adapter contract, a single benchmark schema,
and a single CertOS lifecycle.

\subsection{Adapter contract}
Every supported domain enters ORIUS through the same narrow boundary. The adapter must
parse telemetry, expose reliability-conditioned uncertainty semantics, tighten the
admissible action set, implement repair and fallback, and emit a certificate payload
compatible with the runtime ledger. This is intentionally smaller than a full controller
API. ORIUS is not trying to normalize planning logic; it is trying to normalize safety
interposition.

\subsection{Universal benchmark surface}
The benchmark contract scores safety on the true state and keeps the runtime outputs
comparable across domains. The canonical step-level fields are
\texttt{true\_constraint\_violated}, \texttt{observed\_constraint\_satisfied},
\texttt{true\_margin}, \texttt{observed\_margin}, \texttt{intervened},
\texttt{fallback\_used}, \texttt{certificate\_valid}, \texttt{latency\_us}, and
\texttt{domain\_metrics}. This schema prevents a domain from claiming safety on an
incompatible metric surface.

\subsection{CertOS governance surface}
CertOS makes fallback, certificate validity, lifecycle transitions, and audit continuity
part of the runtime semantics. A domain may choose different repair maps or fallback
modes, but it does not get to opt out of certificate integrity, audit logging, or
release-state discipline. This is why ORIUS reads as a systems paper as well as a
safety paper.

\begin{table*}[t]
\centering
\caption{The three canonical ORIUS surfaces used across all domains.}
\label{tab:ieee-contract-surfaces}
\small
\begin{tabularx}{\textwidth}{@{}p{0.17\textwidth}p{0.23\textwidth}p{0.52\textwidth}@{}}
\toprule
\textbf{Surface} & \textbf{Canonical objects} & \textbf{Purpose in the universal claim}\\
\midrule
Runtime adapter & Telemetry ingest, reliability semantics, uncertainty set, safe action set, repair, fallback, certificate payload & Isolates domain-specific safety geometry while keeping the ORIUS kernel itself domain-neutral.\\
Benchmark schema & True-state violation, observed-state satisfaction, margins, intervention, fallback, certificate validity, latency, domain metrics & Forces all rows to score safety on the true state rather than on controller legality alone.\\
CertOS governance & Lifecycle state, certificate integrity, fallback semantics, audit continuity, release discipline & Prevents runtime safety from collapsing into an undocumented filter or ad hoc emergency stop.\\
\bottomrule
\end{tabularx}
\end{table*}

These interfaces also explain the paper's claim boundary. ORIUS can be universal in
contract even when domains remain unequal in proof depth or data maturity. The parity
gate is therefore not a contradiction of the architecture; it is the mechanism that
stops architectural reuse from being misread as empirical symmetry.
